% -------------------------------------------------------------------------
% ÜBERHOLT (31.08.2026). Maßgeblich ist die eingereichte main.tex. Zahlen,
% Zeilennummern und Angaben zum Overleaf-Stand in dieser Datei stammen aus
% der Zeit VOR der Crosswalk-Korrektur vom 29.08.2026 (35 statt 36 Stadt-
% teile, vollständiger Neulauf) und wurden bewusst NICHT nachgezogen: die
% Datei ist als datierter Entwurf Teil der Arbeitsdokumentation.
% -------------------------------------------------------------------------
% ###########################################################################
% KAPITEL 8 -- VOLLSTAENDIGE NOTIZEN, SCHREIBANLEITUNG UND TABELLEN
% Stand 26.08.2026
%
% EINBAU: ersetzt in main.tex den Bereich von \section{Diskussion} bis
% unmittelbar VOR \section{Fazit und Ausblick} (Stand 26.08.: Zeilen
% 3785-4044). Die beiden Tabellen sind fertiger, kompilierbarer Code; alles
% Uebrige steht als Kommentar und kostet keine Seite.
%
% EINE STELLE AUSSERHALB DIESES BLOCKS IST ANZUFASSEN -- sonst kompiliert die
% Bildunterschrift von Tabelle 8.1 mit einem ?? und einer Warnung:
% Unter \subsection{Forschungsluecke und abgeleitete Erwartungen} (main.tex
% Z. 845) fehlt bisher als einziger Unterabschnitt ein Label. Direkt darunter
% einfuegen:
%     \label{sec:forschungsluecke}
%
% Alle Zahlen in dieser Datei stammen aus results/ des berichteten Laufs
% (#52) und sind am 26.08.2026 gegen die CSV-Dateien geprueft. Wo eine Zahl
% NICHT aus results/ stammt, steht es dabei.
% ###########################################################################


% ===========================================================================
% BLOCK 0 -- DAS SEITENBUDGET. VERBINDLICH.
% ===========================================================================
% Kapitel 8 darf hoechstens 2,5 Seiten belegen. Die folgende Aufteilung ist
% NICHT geschaetzt, sondern am 26.08.2026 mit der Seitengeometrie dieses
% Dokuments nachgemessen (geometry 4/2/4/2 cm, 12pt, onehalfspacing):
%
%   Satzspiegelhoehe        674,3 pt       Zeilenabstand 18,0 pt
%   -> 37,5 Zeilen je Seite, rund 550 Woerter je Seite Fliesstext
%
%   Tabelle 8.1 mit Unterschrift und Float-Abstand      0,25 S.  (gemessen)
%   Tabelle 8.2 mit Unterschrift und Float-Abstand      0,54 S.  (gemessen)
%   1 Ueberschrift Ebene 1 und 3 Ebene 2                0,20 S.
%                                                     ---------
%   fuer Fliesstext bleiben                             1,51 S. = 850 Woerter
%
% VERTEILUNG DER 850 WOERTER -- das ist das eigentliche Budget:
%
%   8.1  3 Absaetze   320 Woerter   (rund 105 je Absatz)
%   8.2  3 Absaetze   280 Woerter   (rund  95 je Absatz)
%   8.3  2 Absaetze   250 Woerter   (rund 125 je Absatz)
%
% Das Schreibrezept verlangt 70 bis 120 Woerter je Absatz -- die Aufteilung
% liegt genau in diesem Fenster. Wer einen Absatz laenger schreibt, nimmt
% ihn einem anderen weg; einen neunten Absatz gibt es nicht.
%
% DREI REGELN, DIE DAS BUDGET TRAGEN -- ohne sie wird es nichts:
%
% 1. JEDE ZAHL STEHT GENAU EINMAL. Entweder in einer Tabelle oder im
%    Fliesstext, nie in beiden. Der Text nennt die RICHTUNG und den GRUND,
%    die Tabelle nennt den WERT. Falsch: "Der Random Forest generalisiert am
%    schlechtesten (R2 0,412 gegen 0,984 im Training)." Richtig: "Der Random
%    Forest generalisiert am schlechtesten und rechnet am laengsten
%    (Tabelle~\ref{tab:erwartung})."
%
% 2. KEINE WIEDERHOLUNG AUS KAPITEL 7. Kapitel 7 berichtet die Ergebnisse,
%    Kapitel 8 deutet sie. Jeder Satz, der eine Zahl aus Kapitel 7 erneut
%    nennt, ohne eine neue Folgerung daraus zu ziehen, wird gestrichen.
%
% 3. LIMITATIONEN GEHEN IN DIE TABELLE, NICHT IN DEN FLIESSTEXT. Der
%    Fliesstext von 8.3 hat zwei Aufgaben und sonst keine: die Groesse der
%    effektiven Stichprobe einordnen und die gemessene Fairnesspruefung
%    berichten. Alles andere ist eine Tabellenzeile.
%
% WENN ES AM ENDE DOCH ZU LANG IST, in dieser Reihenfolge kuerzen:
%   1. Tabelle 8.2 von 13 auf 10 Zeilen -- zuerst fallen "Berichtsumfang",
%      "Schluss auf Personen" und "Gebietszuschnitt" (spart 0,12 S.).
%   2. 8.2 von drei auf zwei Absaetze -- die drei Mechanismen auf zwei
%      Halbsaetze verdichten (spart 0,17 S.).
%   3. Erst danach an 8.1 oder am Fairnessabsatz kuerzen. Beide tragen
%      Bewertungspunkte, die sonst nirgends anfallen.
%
% BLINDABSATZ: zwischen \section{Diskussion} und \subsection 8.1 steht kein
% Text (Vorgabe Schroeter).
%
% FORMALIA (Schroeter 10.08.): keine Kursivschrift, keine
% Anfuehrungszeichen, keine Unterstreichungen, keine Zwischenfazits.
% ===========================================================================


% ===================================================
\section{Diskussion}
\label{sec:diskussion}


% ===========================================================================
% SCHREIBANLEITUNG 8.1 -- VERGLEICHENDE BEWERTUNG DER ALGORITHMEN
% 3 Absaetze, zusammen hoechstens 320 Woerter, danach Tabelle 8.1.
% Beantwortet: Was war zu erwarten, was ist eingetreten, woran liegt es?
% ===========================================================================
%
% ABSATZ 1 -- DIE WIDERLEGTE ERWARTUNG. Gehoert an den ANFANG, nicht ans
% Ende. Abschnitt 2.7 leitet aus Grinsztajn et al. (2022) ab, dass
% baumbasierte Verfahren auf tabellarischen Daten dieser Groessenordnung am
% staerksten abschneiden. Eingetreten ist das Gegenteil: Im Mengenstrang
% schlaegt kein Verfahren die Referenz der Stufe 2, und im Strukturstrang
% kehrt sich das Ergebnis der Kreuzvalidierung auf dem Hold-out um. Das ist
% der interessanteste Befund der Arbeit.
% Erster Satz nennt den Befund, nicht die Erwartung.
%
% ABSATZ 2 -- DIE ERKLAERUNG, UND SIE IST NICHT DIE VERFAHRENSWAHL.
% Das ist die Antwort auf UF4 und das staerkste Argument des Kapitels:
% Zwischen den drei Verfahren liegen 2,2 RMSE, zwischen zwei Spezifikationen
% desselben Verfahrens 19 bis 29. Der Grund ist die multiplikative
% Groessenstruktur: Baumverfahren geben je Blatt einen festen Wert aus und
% koennen "Einsaetze = Bevoelkerung mal Risiko" nicht abbilden; Verfahren mit
% Log-Verknuepfung bekommen die Multiplikation geschenkt.
% ACHTUNG, UEBERHOLTE FASSUNG: Frueher stand hier von einer
% "Spezifikationsasymmetrie" zulasten der Vergleichsverfahren. Seit #43
% modellieren ALLE Verfahren die Rate und multiplizieren zurueck -- die
% Asymmetrie ist konstruktiv beseitigt, nicht nur benannt. NICHT mehr so
% schreiben.
% ZWEITE UEBERHOLTE FASSUNG, ebenfalls nicht mehr schreiben: Der frueher
% vermerkte "schwaechere Rest" der Asymmetrie ist gemessen und existiert
% nicht. B-19: Der Offsetvorteil der Referenz betraegt ueber alle 50 Laeufe
% gepaart -0,0017 RMSE bei der Anzahl und -0,0000 bei der Rate; das negative
% Vorzeichen heisst, die Variante OHNE Offset ist minimal besser. Grund:
% log_bevoelkerung steht ohnehin in PRAEDIKTOREN, der Offset legt nur einen
% Ausgangspunkt fest, den ein frei geschaetzter Koeffizient verschiebt.
% Wenn der Punkt ueberhaupt vorkommt, dann als Fussnote mit dieser Zahl.
%
% ABSATZ 3 -- WARUM FLEXIBILITAET HIER NICHT ZAHLT, UND WAS DAS KOSTET.
% Zwei Haelften, beide gehoeren hierher:
% (a) Die Gegenprobe zur Spezifikation (B-41). In-sample ist Nichtlinearitaet
%     klar nachweisbar; ergaenzt man genau die fehlenden Terme und wertet im
%     selben Rahmen aus, faellt die Prognose um den Faktor drei bis fuenf:
%     linear (12 Terme) RMSE 33,98 -- + Quadrate (22) 101,11 --
%     + Interaktionen (57) 121,63 -- + beides (67) 180,86. Alle 200
%     Anpassungen konvergiert, also kein numerisches Artefakt.
%     Ordnet man die Verfahren nach ungebremster Flexibilitaet, faellt R2
%     monoton: GLM 0,542 -- XGBoost 0,532 -- Ridge 0,511 -- Random Forest
%     0,402 -- Poisson mit Interaktionen, unbestraft, -6,25. Was die
%     Ensembles rettet, ist ihre Regularisierung, nicht ihre
%     Ausdrucksfaehigkeit. Bei 29 unabhaengigen Einheiten ist das das
%     Bias-Varianz-Dilemma in Reinform.
%     DAS IST DER UEBERTRAGBARE METHODISCHE BEITRAG DER ARBEIT: In-sample
%     nachweisbare Nichtlinearitaet ist KEIN hinreichender Grund fuer
%     flexible Verfahren, wenn die Diagnostik geklumpte Daten als unabhaengig
%     behandelt. Ein Satz genuegt, aber er muss stehen.
% (b) Der Preis in Rechenzeit. Der Random Forest braucht 861-mal laenger als
%     Ridge und generalisiert dabei am schlechtesten -- Groessenordnungen
%     statt Prozentpunkte, und damit die belastbarste Aussage des Kapitels.
%     Die Zahlen stehen in der Tabelle, nicht im Satz.
%
% WAS AUS PLATZGRUENDEN NICHT IN DEN FLIESSTEXT KOMMT, ABER IN DIE TABELLE:
% - Unterscheidbarkeit (R-1): keine einzige Verfahrenspaarung ist trennbar.
%   Formulierungsauflage B7 -- NICHT "nicht gezeigt", sondern "nicht
%   gemessen". Sekundaertests Menge: alle sechs nach Holm nicht signifikant;
%   Struktur: RF gegen XGBoost -0,0045, 4 von 10 gewonnen, p 0,625.
% - Ueberanpassung (#51, B6): die R2-Paare Training gegen Kreuzvalidierung.
%   Zu benennende Grenze: Ein Wald mit min_samples_leaf = 1 interpoliert
%   seine Trainingsdaten konstruktionsbedingt; ein hohes Trainings-R2 belegt
%   dort fuer sich genommen keine krankhafte Ueberanpassung. Der Abstand ist
%   zwischen KONFIGURATIONEN und gegen die linearen Modelle zu lesen, nicht
%   als Rangfolge zwischen den Verfahren.
% ===========================================================================

\subsection{Vergleichende Bewertung der Algorithmen}
\label{sec:diskussion-verfahren}

% [ABSATZ 1 -- widerlegte Erwartung]
% [ABSATZ 2 -- Spezifikation schlaegt Verfahrenswahl]
% [ABSATZ 3 -- Flexibilitaet zahlt sich nicht aus, und sie kostet]

\begin{table}[H]
\centering
\caption{Erwartung aus Abschnitt~\ref{sec:forschungsluecke} gegen gemessenen Befund; Gütemaße als Mittel über die zehn Wiederholungen, Zielgröße Anzahl Einsätze}
\label{tab:erwartung}
\footnotesize
\singlespacing
\setlength{\tabcolsep}{4pt}
\begin{tabularx}{\textwidth}{@{}>{\raggedright\arraybackslash}p{2.4cm} >{\raggedright\arraybackslash}p{3.5cm} >{\raggedright\arraybackslash}X@{}}
\toprule
\textbf{Frage} & \textbf{Erwartung} & \textbf{Befund} \\
\midrule
Menge gegen Referenz & Ensembles übertreffen sie & kein Verfahren übertrifft sie: Referenz 33,98 gegen 35,61 bis 37,77 RMSE; Ridge und XGBoost sind signifikant schlechter, der Random Forest nicht unterscheidbar \\
\addlinespace
Struktur gegen Referenz & Ensembles übertreffen sie & in der Kreuzvalidierung ja, $+0{,}031$ und $+0{,}035$ Macro-F1 bei 10 von 10 gewonnenen Wiederholungen --- auf dem Hold-out umgekehrt, $-0{,}070$ und $-0{,}067$ \\
\addlinespace
Verfahren untereinander & Rangfolge erkennbar & keine Paarung trennbar; die Streuungsbereiche überlappen in beiden Strängen \\
\addlinespace
Aufwand & Ensembles teurer & bestätigt, um Größenordnungen: 0,009 s je Anpassung bei Ridge, 0,98 s bei XGBoost, 7,92 s beim Random Forest \\
\addlinespace
Überanpassung & Ensembles anfälliger & bestätigt: $R^2$ Training gegen Kreuzvalidierung 0,826 zu 0,514 bei Ridge, 0,983 zu 0,510 bei XGBoost, 0,984 zu 0,412 beim Random Forest \\
\bottomrule
\end{tabularx}
\end{table}


% ===========================================================================
% SCHREIBANLEITUNG 8.2 -- ERKLAERUNGSBEITRAG DER UNTERSUCHTEN FAKTOREN
% 3 Absaetze, zusammen hoechstens 280 Woerter, KEINE eigene Tabelle.
% Beantwortet UF1. Bis 26.08.2026 hatte dieser Abschnitt als einziger
% ueberhaupt keine Schreibanleitung.
% ===========================================================================
%
% SCHNITT GEGEN 7.4 -- OHNE IHN DOPPELN SICH DIE BEIDEN ABSCHNITTE:
%   7.4 zeigt die Zahlen: SHAP-Attribution und Ablation je Faktorgruppe,
%       mit der Tabelle. 8.2 nennt KEINE dieser Zahlen erneut.
%   8.2 zieht die Folgerung: was die Zahlen ueber die Faktoren aussagen,
%       woran ihre Uneinheitlichkeit liegt und was sie NICHT hergeben.
% Ein Rueckverweis auf die Tabelle in 7.4 ersetzt jede Wiederholung.
%
% ABSATZ 1 -- DIE ANTWORT AUF UF1, UND SIE IST NICHT DIE NAHELIEGENDE.
% FALSCH: "Nur der Kriminalitaetsindex traegt." Im Strukturstrang widerlegt,
%         dort ist er praktisch wertlos (+0,0004 bis +0,0020 Macro-F1).
% RICHTIG: Die drei Faktorgruppen des Exposes tragen dieselbe Information
%          MEHRFACH; welche davon nuetzt, haengt an der zeitlichen Aufloesung
%          und an der Modellform.
% Praezisierung, ohne die die Aussage angreifbar ist: Im Querschnitt
% korrelieren mehrere Merkmale aehnlich stark mit der Zielgroesse
% (anteil_risikogewerbe_pct +0,72, anteil_wohngebaeude_pct -0,79,
% leerstandsquote_pct +0,72 gegenueber +0,66 beim Kriminalitaetsindex). Der
% korrekte Satz lautet: GEGEBEN den Kriminalitaetsindex tragen die uebrigen
% Bloecke nichts Eigenstaendiges mehr bei -- nicht "nur Kriminalitaet haengt
% zusammen".
%
% ABSATZ 2 -- DIE DREI MECHANISMEN, die die Uneinheitlichkeit vollstaendig
% erklaeren. Verdichtet auf drei Halbsaetze, die Zahlen stehen in 7.4:
%   (1) ZEITLICHE AUFLOESUNG. Eindeutige Werte je Stadtteil ueber 132 Monate:
%       Kriminalitaetsindex 128,6 -- ACS-Merkmale und log_bevoelkerung je 4
%       (fuenf Jahrgaenge mit Publikationsversatz) -- die drei baulichen
%       Merkmale 1, also KONSTANT (Land-Use-Snapshot 2020). Der
%       Kriminalitaetsindex ist das einzige monatlich variierende Merkmal.
%   (2) KOLLINEARITAET. Er korreliert mit anteil_risikogewerbe_pct +0,739,
%       leerstandsquote_pct +0,636, armutsquote_pct +0,622,
%       anteil_wohngebaeude_pct -0,603 und ist damit ein SAMMELINDIKATOR, der
%       die soziooekonomische und bauliche Information mittraegt, nur besser
%       aufgeloest. Bei 23 Trainingsstadtteilen kosten redundante Praediktoren
%       mehr Varianz, als sie Verzerrung abbauen (hoechster VIF 12,29).
%   (3) DER OFFSET. Die Groessenkontrolle schadet in der Menge und ist in der
%       Struktur das wertvollste Merkmal. Das Poisson-GLM hat
%       log(Bevoelkerung) ohnehin als Offset -- der Praediktor ist dort ein
%       redundanter freier Parameter. Die Klassifikationsmodelle haben keinen
%       Offset und brauchen die Groesse als Signal.
%
% ABSATZ 3 -- DIE EINSCHRAENKUNG, DIE DIE ANTWORT TRAEGT (R-18).
% Das ist der Absatz, der bisher NIRGENDS vorgesehen war, und er ist der
% wichtigste des Abschnitts: log_kriminalitaetsindex ist das einzige Merkmal,
% dessen Weglassen die Prognose messbar verschlechtert (+24,3 RMSE, 10 von 10
% Wiederholungen). Der gesamte inhaltliche Befund zu UF1 haengt an dieser
% einen Spalte -- und sie stammt aus SFPD Incident Reports, also aus
% polizeilich REGISTRIERTEN Delikten. Registrierte Kriminalitaet ist das
% Produkt aus Tatgeschehen, Anzeigebereitschaft und Kontrollintensitaet; wo
% mehr gestreift wird, wird mehr registriert.
% FORMULIERUNGSAUFLAGE, verbindlich:
%   NICHT: "Kriminalitaet erklaert Feuerwehreinsaetze."
%   SONDERN: "Die registrierte Kriminalitaetsbelastung eines Stadtteils ist
%             der staerkste verfuegbare Indikator." Schwaecher formuliert,
%             aber haltbar.
% Belegbar mit Lum & Isaac (2016), Significance 13(5), 14-19 -- kanonischer
% Nachweis der Rueckkopplung zwischen Polizeieinsatz und registrierten
% Delikten. QUELLE IST NOCH NICHT GEPRUEFT, siehe Block Z.
%
% EIN ZWEITER BEFUND, der in einen Halbsatz passt und viel traegt:
% Die BAULICHEN Merkmale schaden in ALLEN VIER Modellen und BEIDEN Straengen.
% Es sind genau die drei, die aus dem Snapshot eines einzelnen Jahres stammen
% und je Stadtteil konstant sind -- drei Koeffizienten, gefittet auf 23
% Stadtteile Zwischenvarianz. Gleichzeitig eine Limitation der Datenlage und
% ein methodischer Befund. Die Zeile dazu steht in Tabelle 8.2.
%
% ZWEI SAETZE, DIE NICHT FEHLEN DUERFEN (ggf. als Fussnote):
% - Kein Signifikanztest auf der Ablation. Die Testfamilien sind mit #38
%   festgelegt; weitere Tests beruehrten die Holm-Struktur. Berichtet werden
%   Mittelwert, Streuung ueber die zehn Wiederholungsmittel und die Zahl der
%   Wiederholungen mit Verschlechterung.
% - Aus der Ablation folgt NICHT, den Merkmalssatz zu kuerzen. Er kommt aus
%   dem Expose und ist durch die Fairness-Regel gebunden; nachtraeglich zu
%   kuerzen waere eine ergebnisgetriebene Spezifikationswahl.
% ===========================================================================

\subsection{Erklärungsbeitrag der untersuchten Faktoren}
\label{sec:diskussion-faktoren}

% [ABSATZ 1 -- die Faktorgruppen tragen dieselbe Information mehrfach]
% [ABSATZ 2 -- die drei Mechanismen: Aufloesung, Kollinearitaet, Offset]
% [ABSATZ 3 -- registrierte Kriminalitaet, nicht Kriminalitaet]


% ===========================================================================
% SCHREIBANLEITUNG 8.3 -- LIMITATIONEN UND UEBERTRAGBARKEIT
% 2 Absaetze, zusammen hoechstens 250 Woerter, danach Tabelle 8.2.
% Der Fliesstext hat GENAU ZWEI Aufgaben. Alles andere ist Tabellenzeile.
% ===========================================================================
%
% ABSATZ 1 -- DIE EFFEKTIVE STICHPROBE. Die Zahl lautet: rund 38 effektive
% Einheiten bei 4.620 Zeilen (ICC 0,926, Designeffekt 122). Ohne Saison und
% Kriminalitaetsindex nehmen die Merkmale nur 140 verschiedene Werte an.
% (Diese Zahlen stehen NICHT in results/ -- sie sind eine Varianzzerlegung
% der Praediktoren, B-49. Reproduktionsweg dort.)
%
% DER TON ENTSCHEIDET HIER MEHR ALS DIE ZAHL. Drei Fassungen:
%   FALSCH, defensiv: "Leider standen nur 35 Stadtteile zur Verfuegung."
%       Liest sich als Entschuldigung und laedt zur Nachfrage ein.
%   FALSCH, offensiv: "Viele Arbeiten mit 35 Gebietseinheiten berichten
%       p-Werte auf Basis der Zeilenzahl und merken es nicht."
%       Ein Seitenhieb auf die Literatur OHNE BELEG. Angreifbar, nicht
%       souveraen. NICHT SO SCHREIBEN.
%   RICHTIG, sachlich: Die Groessenordnung als Eigenschaft des Gegenstands
%       benennen, die eigene Behandlung dagegenstellen, beides belegen.
%
% GERUEST, KEINE VORLAGE -- eigener Satzbau, eigene Reihenfolge:
%   "Flaechenbezogene Untersuchungen arbeiten typischerweise mit wenigen
%   Dutzend Gebietseinheiten; die effektive Stichprobe liegt hier bei rund
%   38. Entscheidend ist weniger diese Groesse als der Umgang mit ihr:
%   Getestet wird auf den zehn Wiederholungsmitteln statt auf den 50
%   Einzellaeufen, weil letztere Pseudoreplikation waeren; validiert wird
%   stadtteilweise statt zufaellig, weil zufaelliges k-Fold bei
%   Abhaengigkeitsstrukturen zu optimistische Schaetzungen liefert. Die
%   berichtete Streuung quantifiziert die Abhaengigkeit von der
%   Fold-Zuteilung und ist damit eine UNTERGRENZE der Unsicherheit ueber die
%   Grundgesamtheit der Stadtteile."
%
% ERHALTEN BLEIBEN MUESSEN nur die pruefbaren Behauptungen:
%   - die Zahlen: n_eff rund 38, ICC 0,926, Designeffekt 122, 140
%     Merkmalsvektoren
%   - Hurlbert (1984) an der Stelle "warum n = 10 und nicht 50"
%   - Roberts et al. (2017) an der Stelle "warum stadtteilweise, nicht
%     zufaellig" -- diese Quelle IST geprueft und steht bereits in 6.1
%     zitiert, Seite 925
%   - die Richtung: die berichtete Streuung ist eine UNTERGRENZE
%   - kein unbelegter Seitenhieb auf andere Arbeiten
%
% FALLE BEIM UMFORMULIEREN: Aus "Untergrenze der Unsicherheit" wird leicht
% "das Konfidenzintervall ist konservativ" oder "vorsichtig geschaetzt". Das
% ist das GEGENTEIL und waere falsch. Konservativ hiesse zu breit; hier ist
% es zu SCHMAL, weil alle zehn Wiederholungen dieselben 29 Stadtteile teilen
% (B-50). Sichere Wendungen: "eher zu schmal als zu breit", "die
% tatsaechliche Unsicherheit liegt darueber".
%
% EBENFALLS IN DIESEN ABSATZ, ein Satz, weil er der naheliegenden
% Pruefungsfrage die Spitze nimmt: Eine Ausweitung des Zeitraums wuerde die
% Aussagekraft NICHT erhoehen -- die Beschraenkung liegt bei 35
% Analyseeinheiten, nicht bei der Zahl der Beobachtungen. Ein zusaetzliches
% Jahr braechte 420 Zeilen und null zusaetzliche Stadtteile; unter einem
% Stadtteil-Split besteht die Testmenge aus Stadtteilen.
%
% ------------------------------------------------------------------------
% ABSATZ 2 -- DIE GEMESSENE FAIRNESSPRUEFUNG (R-24, B-52).
% Bis 26.08.2026 war dieser Absatz NIRGENDS vorgesehen. Er gehoert hierher,
% weil das Gutachten nach R6/R9 nicht bewertet, ob Probleme auftreten,
% sondern ob sie methodisch geloest wurden -- und ein gemessener und
% verneinter Verdacht ist staerker als ein ungeprueftes Bedenken.
%
% Der naheliegende Vorwurf: Das Modell ist fuer arme Stadtteile ungenauer,
% benachteiligt sie also. Geprueft am 17.08.2026 ueber alle 50 Laeufe, indem
% je Fold das Sozialprofil der Teststadtteile gegen die Fold-Guete gestellt
% wurde.
%   ABSOLUT besteht der Zusammenhang: Spearman-Rho zwischen Armutsquote und
%   RMSE +0,35 bis +0,60 (p < 0,015), ueber alle Verfahren und beide
%   Zielgroessen, die Poisson-Referenz eingeschlossen.
%   RELATIV ZUM NIVEAU verschwindet er vollstaendig: setzt man den RMSE ins
%   Verhaeltnis zur mittleren Einsatzzahl des Folds, liegt Rho zwischen
%   -0,006 und +0,275, und kein Wert ist auf 5 Prozent signifikant
%   (kleinstes p = 0,053).
%   LESART: Der absolute Fehler ist in aermeren Stadtteilen groesser, weil
%   dort mehr Einsaetze stattfinden. Die relative Genauigkeit ist ueber das
%   Sozialgefaelle hinweg gleich. Es liegt kein Hinweis auf systematische
%   Benachteiligung vor.
%
% ZWEI EHRLICHKEITEN, die dazugehoeren und nicht wegfallen duerfen:
%   1. Der relative Fehler ist auf allen Stufen hoch -- im Mittel 0,48 bei
%      der Anzahl und 0,66 bis 0,69 bei der Rate. Das Modell ist nirgends
%      praezise, nicht nur in armen Stadtteilen.
%   2. Es ist eine Pruefung auf FehlerGLEICHHEIT, nicht auf
%      Verteilungsgerechtigkeit eines hypothetischen Einsatzes. Die waere
%      eine andere Frage und wird hier nicht beantwortet.
%
% WICHTIG: Diese Zahlen liegen NICHT in results/. Reproduktionsweg in B-52:
% je (Wiederholung, Fold) die mittlere Armutsquote der Teststadtteile
% berechnen und gegen die Fold-Guete stellen (Spearman, 50 Laeufe).
% Vor der Verwendung im Text einmal nachrechnen und die Zahlen ablegen.
%
% ANSCHLUSS IN EINEM HALBSATZ, wenn der Platz reicht (R-20): Die Abgrenzung
% der Deployment-Phase ist bisher nur eine Umfangsentscheidung; sie ist
% zusaetzlich als ethische Grenze zu formulieren. Das unterscheidet die
% Arbeit von der Predictive-Policing-Kritik, statt ihr auszuweichen: Ein
% zusaetzlicher Loeschzug NUETZT einem Stadtteil, zusaetzliche Streifen
% BELASTEN ihn. Zuteilung nach vorhergesagtem Bedarf ist in der
% Notfallversorgung nicht per se ungerecht, sondern das Ziel.
% Wenn der Platz nicht reicht: in den Ausblick, Kapitel 9.
% ===========================================================================

\subsection{Limitationen und Übertragbarkeit}
\label{sec:limitationen}

% [ABSATZ 1 -- effektive Stichprobe, und warum mehr Zeit nicht hilft]
% [ABSATZ 2 -- Fairness gemessen und verneint, mit den zwei Ehrlichkeiten]

\begin{table}[H]
\centering
\caption{Grenzen der Untersuchung und ihre Wirkung auf die Aussagekraft}
\label{tab:limitationen}
\footnotesize
\singlespacing
\setlength{\tabcolsep}{4pt}
\begin{tabularx}{\textwidth}{@{}>{\raggedright\arraybackslash}p{3.1cm} >{\raggedright\arraybackslash}p{4.4cm} >{\raggedright\arraybackslash}X@{}}
\toprule
\textbf{Grenze} & \textbf{Beleg} & \textbf{Wirkung auf die Aussage} \\
\midrule
Abhängige Wiederholungen & zehn Wiederholungen teilen dieselben 29 Stadtteile & das Konfidenzintervall ist eher zu schmal als zu breit \\
\addlinespace
Extrapolation & 34,6 \% der Testzeilen außerhalb des gelernten Bereichs, 7,6 \% im Hold-out; 9 von 29 Stadtteilen zu 100 \% & begrenzt die Übertragung auf unbekannte Stadtteile \\
\addlinespace
Hold-out und Kreuzvalidierung nicht vergleichbar & 29 statt 23 Trainingsstadtteile, 7,6 statt 34,6 \% Extrapolation & die beiden Zahlenreihen dürfen nicht nebeneinandergestellt werden \\
\addlinespace
Parameterwahl im Hold-out & die Baumverfahren erben die Parameter eines einzelnen Folds, die Referenz hat keine & Schwäche des Hold-out-Verfahrens, nicht der Verfahren; Wirkung nicht gemessen \\
\addlinespace
Einmaliges Tuning & die Parameter stammen aus der ersten Wiederholung und gelten für alle zehn & schwacher Informationsübertrag, alle Verfahren gleichermaßen betroffen \\
\addlinespace
Ungetunte Referenz der Struktur & mit Strafterm 0,314 statt 0,297 Macro-F1 & der Vorsprung in der Kreuzvalidierung halbiert sich auf $+0{,}014$ und $+0{,}020$ \\
\addlinespace
Seltene Klasse & 33 brandbestimmte Monate in der Entwicklung, 2 bis 13 je Testfold & ein Viertel des Macro-F1 ruht auf wenigen Fällen; der Random Forest erreicht auf dem Hold-out bei dieser Klasse 0,173 AUROC \\
\addlinespace
Registrierte Kriminalität & stärkstes Merkmal, $+24{,}3$ RMSE beim Weglassen; Quelle sind Polizeimeldungen & die Aussage gilt für die registrierte Belastung, nicht für Kriminalität \\
\addlinespace
Zeitliche Auflösung der Merkmale & ACS jährlich mit Publikationsversatz, Flächennutzung als Momentaufnahme eines Jahres & die baulichen Merkmale verschlechtern alle vier Modelle in beiden Strängen \\
\addlinespace
Gebietszuschnitt & 35 Analysis Neighborhoods statt 242 Census Tracts & eine andere Einteilung kann andere Koeffizienten und eine andere Rangfolge liefern \\
\addlinespace
Schluss auf Personen & Aggregate je Stadtteil und Monat & keine Aussage über einzelne Haushalte ableitbar \\
\addlinespace
Berichtsumfang & der Mengenstrang ist auf die Anzahl verdichtet, die Rate erscheint als Robustheitsbefund & geprüft: alle Tests auf der Rate stimmen in Richtung und Signifikanzmuster überein \\
\addlinespace
Abweichungen vom Exposé & Stadtteil-Split statt zeitreihengerechter Kreuzvalidierung, zwei statt drei Verfahren in der Klassifikation & begründet in Abschnitt~\ref{sec:rahmen-baselines} und hier; die zeitreihengerechte Variante wird nicht gerechnet \\
\bottomrule
\end{tabularx}
\end{table}


% ###########################################################################
% BLOCK Z -- WAS NICHT IN KAPITEL 8 GEHOERT, WAS OFFEN IST, WAS FEHLT
% Kein Text, nur Merkposten. Vor der Abgabe abarbeiten.
% ###########################################################################
%
% ---------------------------------------------------------------------------
% Z1. AUSDRUECKLICHE GEGENANWEISUNGEN -- NICHT NACH KAPITEL 8
% ---------------------------------------------------------------------------
% - DIE DECKEN (B-48). Gehoeren VOR die Ergebnistabelle in 7.2, nicht in die
%   Limitationen -- sonst lesen sie sich als nachtraegliche Entschuldigung
%   statt als Massstab. Werte: Mehrheitsklasse 0,2233 -- Decke B
%   (Stadtteilwissen) 0,4572 -- Decke A (Label-Rauschen) 0,6404 -- fehlerfrei
%   1,0. Hold-out: Decke A 0,6785, Decke B 0,4219.
% - B-45 (weitere Suchraeume verschlechtern die Prognose) traegt den Vermerk
%   NUR INTERN, geht NICHT in die Arbeit. Nach #52 wird genau ein Lauf
%   berichtet, kein Vorher-Nachher. Uebernehmbar ist allein die allgemeine
%   Lehre, und die ist B-41 (steht in 8.1).
% - Der Seitenhieb auf die Literatur zur effektiven Stichprobe: ohne Beleg,
%   ausdruecklich nicht zu schreiben.
% - Die Extrapolation erklaert KEINEN Verfahrensunterschied (B-31, gemessen
%   und widerlegt: der Zusammenhang zwischen Extrapolationsanteil und RMSE
%   ist bei RIDGE am staerksten, Spearman +0,302, beim Random Forest am
%   schwaechsten, +0,108 bei p 0,456). Sie bleibt eine Eigenschaft des
%   Validierungsrahmens. Diese Erklaerung darf in 8.1 NICHT auftauchen.
%
% ---------------------------------------------------------------------------
% Z2. AUS PLATZGRUENDEN GESTRICHEN -- bewusst, nicht vergessen
% ---------------------------------------------------------------------------
% Diese Punkte sind fuer Kapitel 8 vorgemerkt, passen aber nicht in 2,5
% Seiten. Wenn Platz entsteht, in dieser Reihenfolge aufnehmen:
%   1. Die Anpassungsskala (Alle drei Vergleichsverfahren werden auf der Rate
%      angepasst, das Poisson-GLM auf der Anzahl mit Offset; gemessen wird
%      symmetrisch, angepasst nicht). Zwei Gegenproben liegen vor: im
%      Ratenstrang haben die Baeume die Ausrichtung zwischen Verlust und
%      Metrik und verlieren trotzdem (Hold-out RMSE GLM 0,977 gegen XGBoost
%      1,137, Ridge 1,153, Random Forest 1,561); die Expositions-Ablation
%      kostet den Random Forest 29,2 und XGBoost 19,1 RMSE, die gewaehlte
%      Spezifikation ist also die fuer die Vergleichsverfahren guenstigere.
%      Eine expositionsgewichtete Variante wurde NICHT erhoben -- #43 legt
%      eine Spezifikation fuer alle fest, nachtraeglich weitere zu messen und
%      unter ihnen zu waehlen waere ergebnisgetrieben.
%   2. Die Rueckkopplung bei hypothetischem Einsatz (R-23): mehr Praesenz,
%      mehr registrierte Ereignisse, Bestaetigung der Vorhersage. Beim Brand
%      schwaecher als bei der Polizei, bei Fehlalarmen (79 Prozent, dominante
%      Klasse) denkbar. Laut Register gehoert das ohnehin in den AUSBLICK,
%      Kapitel 9, nicht in Kapitel 8.
%   3. Das fehlende lineare Referenzverfahren der Klassifikation (#31): Die
%      Symmetrie linear/Bagging/Boosting besteht nur im Regressionsstrang;
%      der Vergleich in der Klassifikation ist einer zwischen zwei
%      Baumverfahren, die einander konstruktionsbedingt aehnlicher sind.
%      Steht bereits in 6.2 -- in 8.3 nur, wenn Platz ist.
%   4. Die Reflexion der beiden Dokumentationsfehler (B-43/B-44): eine
%      Information lebte an zwei Orten, nur einer wurde gepflegt. Trifft
%      Gutachten-R9 (Reflexion im Verlauf). ENTSCHEIDUNG OFFEN, ob ein selbst
%      berichteter Prozessfehler ueberhaupt in die Arbeit soll.
%
% ---------------------------------------------------------------------------
% Z3. ZAHLENKONFLIKTE -- VOR DEM SCHREIBEN KLAEREN
% ---------------------------------------------------------------------------
% (a) RESET-TEST, ZEILENZAHL. 06_RISIKEN.md (Z. 392), 07_BEFUNDE.md B-41 und
%     der Kommentar in main.tex Z. 2456 nennen 3.828 Zeilen. Richtig sind
%     3.036 -- die Eignungspruefung lief auf den 23 Trainingsstadtteilen von
%     Fold 1 (results/eignungspruefung/eignungspruefung.md, Z. 3). 3.828 ist
%     die TRAININGSGROESSE DES HOLD-OUTS (29 mal 132, holdout.csv). Der
%     Fliesstext in 6.2 (main.tex Z. 2568-2571) ist bereits richtig; die drei
%     Kommentar- und Notizstellen sind nachzuziehen.
% (b) STADTTEIL-MITTELWERT. #29 nennt R2 0,888, die Praezisierung vom
%     04.08.2026 und 06_RISIKEN.md nennen 0,925. Genau diese Zahl traegt
%     Auflage A. Vor der Verwendung festlegen, welche gilt.
% (c) VORSPRUNG DER ENSEMBLES IN DER STRUKTUR. #48 und 07_BEFUNDE B-29 nennen
%     +0,0304/+0,0371 bzw. +0,0296/+0,0362 -- beides Werte vom 06.08. Der
%     BERICHTETE Lauf (#52) liefert +0,0306 und +0,0350
%     (klassifikation/vergleich.csv). In der Arbeit gelten die letzteren.
% (d) DIE GETUNTE REFERENZ 0,314 stammt aus der Messung vom 06.08. und wurde
%     nach #48 bewusst NICHT nachgerechnet. Sie ist damit eine Zahl aus einem
%     anderen Lauf als die uebrigen. Beim Schreiben so kennzeichnen.
% (e) DISPERSIONSINDEX. annahmen.csv nennt 54,2, #42 nennt 62,8. Es gilt
%     54,2 (berichteter Lauf, Fold-1-Trainingsstadtteile).
%
% ---------------------------------------------------------------------------
% Z4. QUELLEN, DIE DIESES KAPITEL BRAUCHT
% ---------------------------------------------------------------------------
% Roberts et al. (2017) ist geprueft, im Literaturverzeichnis und in 6.1 mit
% Seite 925 zitiert -- die einzige der folgenden, die verwendbar ist.
%
% NICHT GEPRUEFT, kein bib-Eintrag, keine Seitenangabe. Nach der geltenden
% Regel darf keine davon vor der Volltextpruefung in den Text:
%   Hurlbert, S. H. (1984): Pseudoreplication and the Design of Ecological
%     Field Experiments. Ecological Monographs 54(2), 187-211.
%     -> braucht 8.3 Absatz 1 fuer "warum n = 10 und nicht 50"
%   Lum, K. & Isaac, W. (2016): To predict and serve? Significance 13(5),
%     14-19.  -> braucht 8.2 Absatz 3
%   Robinson, W. S. (1950): Ecological Correlations and the Behavior of
%     Individuals. American Sociological Review 15(3), 351-357.
%     -> staerkt die Tabellenzeile "Schluss auf Personen"
%   Openshaw, S. (1984): Ecological Fallacies and the Analysis of Areal
%     Census Data. Environment and Planning A 16(1), 17-31.
%     -> staerkt die Tabellenzeile "Gebietszuschnitt"
%   Nadeau, C. & Bengio, Y. (2003): Inference for the Generalization Error.
%     -> staerkt die Tabellenzeile "Abhaengige Wiederholungen"
%
% ENTLASTUNG: Die beiden Tabellenzeilen zum Gebietszuschnitt und zum Schluss
% auf Personen sind Aussagen ueber das EIGENE Design und tragen auch ohne
% Zitat. Nur Hurlbert und Lum & Isaac werden wirklich gebraucht; ohne sie ist
% der jeweilige Satz abzuschwaechen, nicht zu streichen.
%
% ---------------------------------------------------------------------------
% Z5. SCHREIBREIHENFOLGE UND ABHAENGIGKEITEN
% ---------------------------------------------------------------------------
% - Abschnitt 2.7 (Forschungsluecke) wird ZULETZT geschrieben, nach 8.1. Die
%   dort abgeleiteten Erwartungen sind genau die, die 8.1 widerlegt -- sie
%   muessen zusammenpassen (09_SCHREIBPLAN.md Z. 74).
% - AUFLAGE A (Schroeter, 04.08.2026, woertlich): "Wichtig ist, dass Sie die
%   Abweichung von der urspruenglich angekuendigten zeitreihengerechten
%   Kreuzvalidierung transparent erlaeutern." Das begruendete Verwerfen ist
%   VERLANGT, nicht optional. Es steht als letzte Zeile in Tabelle 8.2 -- die
%   Zeile allein erfuellt die Auflage nur, wenn der Beleg mit dranhaengt:
%   92,5 Prozent der Varianz liegen zwischen den Stadtteilen, ein Zeitschnitt
%   wuerde im Kern die zeitliche Stabilitaet messen und die Forschungsfrage
%   nicht pruefen. Pruefen, ob dieser Beleg in 5.4 schon steht -- wenn ja,
%   genuegt in 8.3 der Verweis; wenn nein, muss er in die Tabellenzeile.
% - DARSTELLUNGSREGEL (Schroeter, 10.08.2026): keine grosse iterative
%   Entwicklung im Text; Kapitel 6 beschreibt die finale Spezifikation, die
%   Reflexion ist KONZENTRIERT IN KAPITEL 8. Damit gehoeren #42
%   (Verlustfunktionen) und #43 (Exposition) hierher -- beide ausdruecklich
%   NACH einem Modelllauf getroffen und in der Arbeit als solches zu
%   kennzeichnen. Bei 2,5 Seiten reicht dafuer ein Halbsatz in 8.1 Absatz 2.
% - KI-VERZEICHNIS: Wird eine hier vorformulierte Passage als Geruest
%   genutzt, gehoert ein Eintrag in die Tabelle am Ende des Dokuments, nach
%   dem Muster der vorhandenen Eintraege mit Prompt und Ergebnisverweis.
% - GLOBALE UMBENENNUNG (Schroeter, 24.08.2026): Quelltext/Quellcode wird zu
%   "Code 1: xy". Betrifft auch die \ref-Aufrufe in diesem Kapitel, falls
%   welche dazukommen.
%
% ---------------------------------------------------------------------------
% Z6. PRUEFUNG NACH DEM SCHREIBEN
% ---------------------------------------------------------------------------
% 1. Zweimal kompilieren (Verzeichnisse stabilisieren sich erst im zweiten
%    Lauf), dann die Seitenzahl von \section{Diskussion} bis
%    \section{Fazit und Ausblick} nachmessen. Ziel: hoechstens 2,5.
% 2. Beide Tabellen muessen auf ihrer Seite stehen bleiben. Faellt eine um,
%    ist der Float an den ANFANG seines Blocks zu setzen, nicht ans Ende --
%    [H] allein schiebt sie nur auf die naechste Seite.
% 3. Jede Zahl im Fliesstext gegen die Tabellen pruefen: steht sie zweimal,
%    faellt die im Fliesstext weg.
% 4. Gegen Kapitel 7 lesen: Wird eine Ergebniszahl wiederholt, ohne dass eine
%    neue Folgerung daraus entsteht, streichen.
% ###########################################################################
